\section{问题三的模型建立与求解}
\label{sec:question-three}

\subsection{问题分析与模型输入}
问题二确定了每张图在各核数下的分核与核内顺序；问题三增加共享只读Cache，关键变化是同一份张量在不同核心上的读取时刻和读取来源。为分开观察硬件变化与调度调整，沿用同图同核的问题二完整最终计划$X_B^*$，包括操作到子图的映射、子图到核心的映射和各核内顺序。先在C环境评价这份计划，再围绕共享输入产生新候选，由此得到$F_B(X_B^*)$、$F_C(X_B^*)$和最终的$F_C(X_C^*)$。

共享只读Cache容量为1048576 B，读取带宽为250 B/cycle；DDR带宽仍为60 B/cycle，两个读取池独立。Cache初始为空，按先进先出淘汰。逻辑张量标识用于命中查询，实际COPY事件则来自问题二的完整核内展开，Spill产生的物理版本也由题给官方评估器处理。这里的FIFO方程和三段比值属于模型层，有限候选的迁移、合并、调序与排序属于候选求解器，工期、命中和淘汰字段则只取官方C评估器的事件结果。我们逐组保存B最终计划哈希、C初始计划哈希、三次评价的工期及命中/未命中字节，不以相同候选名称代替计划身份核验。
\begin{table}[H]\centering\caption{问题三的固定输入和配对量}\label{tab:c-inputs}
\begin{tabularx}{\linewidth}{lY}\toprule
对象 & 值或判定条件\\\midrule
同图同核起点 & 问题二最终计划$X_B^*$；B最终哈希等于C初始哈希\\
共享Cache & 容量1048576 B；读带宽250 B/cycle；初始队列为空\\
外部存储 & DDR带宽60 B/cycle；写出与未命中读仍走DDR\\
正式响应 & 总工期、结构与恢复搬运、命中和未命中字节、完整计划哈希\\
选择预算 & 常规额度为$\max\{\min(q,2),m+1\}$；$m$为去重后必评锚点数\\\bottomrule
\end{tabularx}\end{table}

\subsection{基于查询与完成事件的FIFO模型}
设$\mathcal C(\tau)$为时刻$\tau$按插入次序排列的逻辑张量队列。合格读入事件$e$在发射时刻$q_e$查询$t_e$，其读取来源由
\begin{equation}
H_e=\mathbf 1\{t_e\in\mathcal C(q_e)\},\qquad
B_e=\begin{cases}250,&H_e=1,\\60,&H_e=0\end{cases}\ \mathrm{B/cycle}
\label{eq:c-query}
\end{equation}
确定。命中读进入Cache池，未命中读进入DDR池；命中只改变本次读取所用的服务池，不把条目移到队尾，查询也不改变FIFO次序。在该COPY\_IN完成时，评估器再尝试把张量插入队尾：若当前仍在队列中则不改变次序；若已经被其他完成事件淘汰，则可能重新插入；若大小超过容量则不插入；空间不足时先逐项淘汰队首。同一时刻的完成、插入和新发射按题给评估器的固定事件次序处理，不能把尚未完成的在途读视为已命中。公式\eqref{eq:c-query}中的$B_e$只是事件选择的服务带宽，不是整图工期的直接除数。

考虑容量为$2b$、三个张量$x,y,z$大小均为$b$的符号例。已有队列$[x,y]$，对$x$的命中不会将其移到队尾。随后$z$完成并插入，队列变为$[y,z]$；如果先前对$x$的在途读取此时才完成，它又能在队尾重新插入，队列变为$[z,x]$。同一组张量的访问次数并不足以确定这段队列演化，完成事件的先后与查询先后都必须考虑。

更关键的是并发查询。两个核心在首个读入完成前向空Cache发射相同张量的请求，两次都未命中；只有后发请求在首次完成插入之后查询且条目未被淘汰，才可能命中。记合格读入总命中字节为$Q_L$，未命中字节为$Q_D$，则报告的字节命中率为
\begin{equation}
h_{\rm byte}=\frac{Q_L}{Q_L+Q_D}.
\label{eq:hit-rate}
\end{equation}
分母为零时遵循官方结果字段的零值约定。这个指标描述读取服务来源，不能从官方“额外搬运”中扣除$Q_L$；搬运量仍按实际生成的COPY事件统计。

DDR和Cache两个共享服务池各自按活动请求分配参考工作量。若池$p$中同时有$m_p(\tau)$个未完成请求，请求$e$的剩余参考工作量$R_e$满足
\begin{equation}
\frac{\mathrm dR_e}{\mathrm d\tau}=-\frac{1}{m_p(\tau)},\qquad
R_e(q_e)=\max\left\{1,\left\lceil\frac{b_e}{B_e}\right\rceil\right\}.
\label{eq:c-service}
\end{equation}
一次命中既可能减轻DDR池竞争，也可能改变后继操作的发射时刻，从而改变关键路径、另一池负载及后续FIFO状态。因此即使$h_{\rm byte}>0$，总工期仍可能不变。

\begin{algorithm}[H]\caption{题给C评估器的只读FIFO事件规则}\label{alg:c-fifo}
\begin{algorithmic}[1]
\REQUIRE 已展开的核级执行图、COPY事件、Cache容量及两类带宽
\ENSURE 完成时间线、Cache事件、命中与未命中字节
\STATE 初始化空FIFO队列、DDR池与Cache读池
\WHILE{仍有操作未完成}
\STATE 按题给评估器的事件次序处理当前完成事件，推进相应池的剩余工作量
\FOR{当前完成的每个合格COPY\_IN}
\IF{逻辑张量不在队列且大小不超过容量}
\STATE 从队首淘汰至空间充足，并插入该张量
\ENDIF
\ENDFOR
\STATE 更新依赖、流水线与跨核释放就绪状态
\FOR{当前可发射的合格COPY\_IN}
\STATE 在发射时查询队列，选择Cache读池或DDR池并累计字节
\ENDFOR
\STATE 重排两池的预计完成事件并推进事件时间
\ENDWHILE
\end{algorithmic}\end{algorithm}
\subsubsection{Cache查询与完成时序}
为把公式\eqref{eq:c-query}与事件先后联系起来，构造一个不参与100图统计的解析小例：三个核心各需读同一份60000 B张量$x$，Cache容量足以容纳$x$。独占DDR读取为$60000/60=1000$ cycle，独占Cache读取为$60000/250=240$ cycle。若三核在时刻0同时发射，三次查询都看到空队列，三个DDR请求平分服务，读取均在时刻3000完成，命中字节为0。若第一核在时刻0发射并于1000完成插入，其余两核在时刻1000才发射且$x$未被淘汰，则两次都命中；两个Cache请求平分读池，各需480 cycle，读取于1480完成。表\ref{tab:c-query-example}把发射、插入和完成三个状态明确分开。
\begin{table}[H]\centering\caption{三核重复读取60000 B张量的解析时序}\label{tab:c-query-example}
\begin{tabularx}{\linewidth}{lYrrr}\toprule
情形 & 发射与插入顺序 & DDR读取 & Cache读取 & 最后一次读完成/cycle\\\midrule
同时查询 & 三核均在0发射，首次插入发生在完成时 & 3次 & 0次 & 3000\\
错峰查询 & 第一核0发射、1000插入，其余两核1000发射 & 1次 & 2次 & 1480\\\bottomrule
\end{tabularx}\end{table}
\begin{figure}[H]\centering
\begin{tikzpicture}[x=1.65cm,y=.9cm,>=Stealth,every node/.style={font=\small}]
\draw[->] (0,0)--(3.5,0) node[right] {$\tau/1000$ cycle};
\foreach \x/\lab in {0/0,1/1,1.48/1.48,3/3}{\draw (\x,.06)--(\x,-.06) node[below] {\lab};}
\node[left] at (0,1.0) {错峰查询};
\node[left] at (0,2.4) {同时查询};
\draw[blue!65!black,very thick] (0,2.4)--(3,2.4);
\node[above] at (1.5,2.4) {3次DDR未命中并发};
\draw[blue!65!black,very thick] (0,1.0)--(1,1.0);
\node[above] at (.5,1.0) {1次DDR};
\draw[green!55!black,very thick] (1,1.0)--(1.48,1.0);
\node[anchor=west] at (1.58,1.0) {2次Cache命中};
\draw[dashed] (1,.12)--(1,2.9);
\node[anchor=west] at (1.65,3.25) {首读完成并插入：$\tau=1000$};
\end{tikzpicture}
\caption{相同访问次数下的两种查询顺序；仅为解析示例}\label{fig:c-query-timeline}
\end{figure}
这里的错峰发射是给定的时序条件，不是调度器可以免费插入的“空等”动作。在真实计划中，后续查询的延迟必须来自前置计算、依赖或队列中的其他操作；这些操作本身也占资源。解析例只证明查询时刻会改变命中与读完成时间，不给出真实图的整体工期。

\subsection{共享输入方案与求解}
\textbf{Step 1：生成候选。}以$X_B^*$为种子，尝试迁移、同核合并、核内调序以及跨核分散，并查找跨核重复读取的逻辑张量。对按张量大小和潜在重复读量排序的有限张量集合，枚举相关消费者组或节点在现有核心序列中的前移、后移和跨核插入位置；对满足容量条件的共享输入，再指定一个领读核，将其消费者移至该核序列前部，并把其他消费者后移到各自序列尾部，形成错峰查询候选。候选只保留合法映射与顺序；共享张量集合随当前Cache容量过滤，因而容量扰动会同步改变候选组合。不合法的构造不进入完整评价。

\textbf{Step 2：轻量排序。}按路径递推、最大核心计算量、结构COPY的DDR参考工作量和容量压力估计候选代价；C另外在收缩后的商图上回放逻辑COPY查询的相对时序，记录查询时的FIFO命中代理、未命中量和潜在Cache服务工作量。该回放保留每核顺序、跨核释放和完成后插入规则，但不替代官方事件池的并发服务；DDR参考工作量仍按结构COPY计算，代理命中不会从结构字节中扣除。正式工期、命中和淘汰由题给问题三事件模拟给出。

\textbf{Step 3：官方评价与选优。}把B最终计划$X_B^*$、整图、分量装桶和低核数继承计划中去重后的方案置于必评锚点，其余C候选按轻量分数排列。设去重后必评数为$m$、请求额度为$q$，常规完整调用额度为$\max\{\min(q,2),m+1\}$。profile失败只记录错误，不占完整调用额度；完整官方调用失败占用一次额度。首个成功结果记为初始方案，最终按工期、额外搬运和固定候选名的顺序在全部成功结果中选定$X_C^*$。首个官方C结果后，从其Cache事件中提取重复未命中张量，生成事件锚定的合法调序候选；若仍有额度，再按候选代表补评。下文500组B/C配对采用同一B计划作为C起点，并逐组核对计划身份；该配对结果用于同计划效应分解，与逐图继承后的总体曲线分开报告。

\begin{algorithm}[H]\caption{问题三候选选择与同计划对照}\label{alg:c-select}
\begin{algorithmic}[1]
\REQUIRE 问题二最终计划$X_B^*$、图、核数、固定硬件参数及请求完整调用额度$q$
\ENSURE 问题三选中计划$X_C^*$及三段工期比
\STATE 保留$X_B^*$，生成迁移、同核合并、调序、分散及领读核错峰代表
\STATE 逐个检查合法性、按完整计划去重，计算轻量排序分数
\STATE 将$X_B^*$、整图、分量装桶和低核数继承计划置于去重后的必评队列
\STATE 令$m$为必评数，完整调用额度为$\max\{\min(q,2),m+1\}$
\WHILE{未达到常规完整调用额度且还有未尝试候选}
\STATE 取必评队列或轻量排序最靠前的未尝试候选，先profile后调用官方C评估器
\STATE profile失败不占额度；完整评价保存成功结果或完整错误记录
\ENDWHILE
\STATE 从首个成功结果提取重复未命中事件，生成事件锚定调序候选并在剩余额度内补评
\STATE 在成功结果中以工期、额外搬运、固定候选名顺序选$X_C^*$
\STATE 核对B最终哈希与C初始哈希，计算式\eqref{eq:cache-factors}
\end{algorithmic}\end{algorithm}
伪代码已列出从B最终计划、候选生成、轻量回放到成功集合选优的完整步骤；其后的分析围绕同计划配对与事件级字段展开。

\subsubsection{共享输入的查询时序}
仍用60000 B共享输入，令五个V流水线操作各计算1000 cycle，暂不计输出。下述同时发射与错峰发射是两种给定的局部请求时序，用来比较服务池的边界工期，并非同一张图上由调度器任意插入等待后得到的可行方案。全部聚到一核，只读一次DDR却顺序完成五个计算，理论完成时间为$1000+5\times1000=6000$ cycle。五核分散且同时查询空Cache时，五次未命中读平分DDR，读阶段持续$5\times60000/60=5000$ cycle，再各算1000 cycle，总时间同为6000。若只有第一核在0查询，其余四核由外部前置关系约束到1000才查询，四次命中合计240000 B，Cache读阶段需$240000/250=960$ cycle，后四核在2960完成。两种发射方式的差异说明，读取时序与Cache服务能力共同决定局部读计算完成时间。
\begin{table}[H]\centering\caption{五个V操作共享60000 B输入的解析边界（非正式实验）}\label{tab:c-five-example}
\begin{tabularx}{\linewidth}{lYr}\toprule
结构 & 读取与执行条件 & 解析工期/cycle\\\midrule
单核聚合 & 一次DDR读取，五次V操作串行 & 6000\\
五核分散、无Cache & 五次DDR读同时开始、共享60 B/cycle & 6000\\
五核分散、Cache同时查询 & 空队列下五次均未命中 & 6000\\
五核分散、其余四核错峰 & 后四次命中，Cache带宽250 B/cycle & 2960\\
同上、带宽减半 & 后四次命中，Cache带宽125 B/cycle & 3920\\
同上、带宽加倍 & 后四次命中，Cache带宽500 B/cycle & 2480\\
同上、容量不足60000 B & 张量不能插入，后四次仍未命中 & 6000\\\bottomrule
\end{tabularx}\end{table}
五核分散在250 B/cycle条件下，第一核于2000完成计算；其余四核从1000起共读960 cycle，再计算1000 cycle，于2960完成。带宽减半和加倍时，将960分别改为1920、480，即得3920、2480。这个局部算例说明为何同时保留聚合和调序候选；真实图中前置操作、输出及其他DDR请求均进入官方事件评价，表内数字不作为其工期预测。

\subsection{同计划对照与调度适配}
同一图同一核数取得三个官方工期：$F_B(X_B^*)$、$F_C(X_B^*)$、$F_C(X_C^*)$。定义
\begin{equation}
\begin{split}
S_{\rm hw}&=\frac{F_B(X_B^*)}{F_C(X_B^*)},\\
S_{\rm adapt}&=\frac{F_C(X_B^*)}{F_C(X_C^*)},\\
S_{\rm opt}&=\frac{F_B(X_B^*)}{F_C(X_C^*)}=S_{\rm hw}S_{\rm adapt}.
\end{split}
\label{eq:cache-factors}
\end{equation}
第一个比值只改硬件场景，第二个只在C中换计划，最后一个是B最终到C最终的综合工期比。500组配对逐组核对计划哈希，乘积恒等式的最大数值残差为$2.22\times10^{-16}$。表\ref{tab:cache-factors}每列分别对100张图的逐例比值取算术平均；\emph{列均值不再满足乘积恒等式}。
\begin{table}[H]\centering\caption{每核100图的Cache收益三段分解}\label{tab:cache-factors}
\begin{tabular}{crrrr}\toprule
核数 & 硬件效应均值 & 适配效应均值 & 综合效应均值 & 适配改善图数\\\midrule
1 & 1.0044 & 1.0249 & 1.0293 & 31\\
2 & 1.0068 & 1.0000 & 1.0068 & 1\\
3 & 1.0088 & 1.0011 & 1.0099 & 1\\
4 & 1.0138 & 1.0011 & 1.0150 & 7\\
5 & 1.0125 & 1.0011 & 1.0136 & 7\\\bottomrule
\end{tabular}\end{table}
\samplefigure{cache_decomposition}{同计划硬件效应、C内调度适配及综合比值}{fig:cache-decomposition}

为避免把比值均值误读为绝对工期，图\ref{fig:cache-direct-comparison}同时给出三种配对量：左图是$F_B(X_B^*)$、$F_C(X_B^*)$和$F_C(X_C^*)$的绝对平均工期，右图是后二者相对$B(X_B^*)$工期的逐图比值均值。左图反映图规模和核数共同造成的周期差异，右图才对应同计划硬件效应与C内调度适配的相对变化。
\samplefigure{cache_direct_comparison}{问题三无Cache、固定计划Cache与重调度Cache的同核直接比较}{fig:cache-direct-comparison}

均值给出了收益幅度，却不能表明有多少图真正缩短了工期。图\ref{fig:cache-coverage-radar}用分组条形图表示每核100图中工期严格缩短的占比，三类效应采用相同的计数口径。五核硬件效应、C内适配与综合改善分别覆盖54、7、58图；单核有31图受益于C内适配，而二、三核各只有1图。硬件与适配的改善图有重合，不能把两类计数直接相加。
\samplefigure{cache_effect_coverage_radar}{五个核数下硬件、适配与综合工期改善的逐图覆盖率分组条形图}{fig:cache-coverage-radar}

五核100图中，$S_{\rm hw}>1$有54图，恰为1有46图，没有退化图；$C(X_B^*)$命中字节加权占比为16.79\%。图\ref{fig:cache-hardware-pie}用分类条形图拆分54图的改善幅度：36图工期降幅不足1\%，改善图工期降幅中位数为0.4428\%，平均工期降幅为2.0074\%；若将同一批逐图工期比倒数后换算为速度增幅，均值为2.31\%。图中的图数构成与工期降幅采用不同坐标量纲，分别读取。
\samplefigure{cache_hw_effect_pie}{五核硬件效应的图数构成及改善幅度分类条形图}{fig:cache-hardware-pie}

对\Case{080}，固定B计划的工期由113455降至90408 cycle，$S_{\rm hw}=1.2549$，对应工期降幅20.3138\%；对\Case{055}，$C(X_B^*)$存在189952字节命中，但工期两侧均为25635 cycle。现有汇总结果支持的结论是：命中不一定位于决定尾部完成时刻的关键路径；事件级瓶颈在图中不另作单一归因。图\ref{fig:cache-hit-gain}同时给出五核100图的$C(X_B^*)$命中率与同计划工期改善，零收益点并未从散点中删去。
\samplefigure{cache_hit_gain_scatter}{五核字节命中率与同计划工期改善的逐图关系}{fig:cache-hit-gain}

\begin{table}[H]\centering\caption{四张代表图的五核三段工期与Cache访问}\label{tab:c-cases}
\begin{tabular}{lrrrrr}\toprule
图 & $F_B(X_B^*)$ & $F_C(X_B^*)$ & $F_C(X_C^*)$ & $C(X_B^*)$命中/B & $S_{\rm hw}$\\\midrule
\Case{026} & 45246 & 44374 & 41655 & 21936 & 1.0197\\
\Case{053} & 1472686 & 1428109 & 1428109 & 5155968 & 1.0312\\
\Case{055} & 25635 & 25635 & 25635 & 189952 & 1.0000\\
\Case{080} & 113455 & 90408 & 90408 & 2228224 & 1.2549\\\bottomrule
\end{tabular}\end{table}
表\ref{tab:c-cases}把影响分成三种可见情形：\Case{055}读来源有变化但完成时间不变；\Case{080}固定计划的硬件效应较大；\Case{026}则在硬件效应之外还有C内调度适配。\Case{053}命中字节较多，硬件效应却仅1.0312，进一步说明命中量不是时间收益的直接比例尺。本表将结论限定到方案级字段能够直接支持的范围。

把同计划比较扩展到全部500个图-核组合，图\ref{fig:cache-hit-gain-3d}按核数分面，在二维坐标中表示$C(X_B^*)$字节命中率与$100[1-F_C(X_B^*)/F_B(X_B^*)]$工期降幅。205组缩短了工期，另有141组虽发生命中但工期不变；154组既无命中也无工期变化，在各分面原点叠印。这500组中，所有同计划工期缩短的结果都伴随非零命中，但命中不足以推出工期一定下降；单凭跨图散点也不能判断某个读取事件是否位于关键路径。
\samplefigure{cache_hit_gain_3d}{500组同计划配对的分核命中率与工期降幅分面散点图}{fig:cache-hit-gain-3d}

\Case{026}则说明第二段效应：同一B计划在C下的起点为44374 cycle，选中重复读取相关候选后为41655 cycle，$S_{\rm adapt}=1.0653$。五核仅7图出现适配改善，说明这类收益集中在少数计划。500组总计有205组硬件效应改善、295组持平；47组适配改善、453组持平；综合改善232组、持平268组。收益的稀疏性也是报告逐例表而非只报总均值的原因。

进一步检查两段收益是否发生在同一批图。以$S_{\rm hw}>1$和$S_{\rm adapt}>1$对五核100图交叉分类，等于1记为未改善。表\ref{tab:c-effect-cross}的综合工期降幅先按每图的$100[1-F_C(X_C^*)/F_B(X_B^*)]$计算，再在组内等权平均。
\begin{table}[H]\centering\small
\caption{五核100图的硬件效应与调度适配交叉分布}
\label{tab:c-effect-cross}
\begin{tabular}{lrr}\toprule
效应分类 & 图数 & 综合工期降幅均值/\%\\\midrule
仅同计划硬件效应 & 51 & 2.0778\\
仅C内调度适配 & 4 & 1.1170\\
两种效应兼有 & 3 & 2.9412\\
两种效应均无 & 42 & 0.0000\\\bottomrule
\end{tabular}
\end{table}
硬件效应改善的54图与适配改善的7图有3图重合，因此综合改善为58图。\Case{040}的B最终与C同计划工期均为110337 cycle，更换C计划后降至110105 cycle，最终计划的命中字节为0。这232 cycle来自已评计划之间的差异，不能直接归为Cache命中收益；51张图则在保留B计划时就已缩短工期。交叉分类表明，是否更换计划与是否受益于Cache硬件是两个独立的判断步骤。

\subsection{问题三结果}
问题三采用与前两问一致的逐图继承口径：对目标核数$K$，在已评价的$1,\ldots,K$核方案及整图单核基准中取Makespan最短者，低核方案以空核列表补齐。相对统一整图单核基准的五核逐例加速比均值为3.4272、中位数4.0719，所有图均不慢于基准。单核C保留真实工期比；A、B单核曲线另按题面显示1。表\ref{tab:c-results}的初始到最终降幅是首个成功官方候选与最终已评最优候选的比较，并非移除某个机制的消融。
\begin{table}[H]\centering\caption{问题三：100图逐核官方结果（慢例按未舍入工期严格比较）}\label{tab:c-results}
\begin{tabular}{crrrr}\toprule
核数 & 平均加速比 & 中位数 & 慢于整图单核图数 & 初始至最终平均工期降幅\\\midrule
1 & 1.0468 & 1.0000 & 0 & 2.1873\%\\
2 & 1.7232 & 1.9703 & 0 & 0.0215\%\\
3 & 2.3335 & 2.8340 & 0 & 0.0215\%\\
4 & 2.9496 & 3.6732 & 0 & 0.0215\%\\
5 & 3.4272 & 4.0719 & 0 & 0.0828\%\\\bottomrule
\end{tabular}\end{table}
\samplefigure{speedup_C}{问题三加速比均值、中位数及均值95\%实例重采样区间}{fig:speedup-c}

图\ref{fig:speedup-c}的浅色带来自固定100张图的1500次有放回实例重采样，反映图集合组成改变时的均值波动，不代表芯片重复运行噪声。均值由4核的2.9496升至5核的3.4272，增加0.4776；由4核到5核有69张图工期缩短、31张图保持不变，没有图因增加一个核心而变慢。

\subsection{问题三灵敏度分析：L2容量与带宽}
固定六张代表图的五核B最终计划，分别把L2 Cache容量与读带宽减半、加倍，考察硬件参数对命中、工期和切图调度的影响。六图按五核B计划的基线工期、初始命中率和规模分层，覆盖小工期低命中图（012、081）、中等工期高命中图（051）、容量边界图（062）以及大工期图（016、067）；样本在实验前固定，未按扰动后的结果挑选。这里“固定计划”只改变硬件参数，“重优化”才允许重新生成C候选；24个固定计划结果全部成功，每个参数组均有成功的重优化候选，共36条成功、12条因官方依赖环检查失败。24组重优化最终全部选回固定B计划，因而本组实验观察到的直接结论是：在这六张代表图和有限候选范围内，容量、带宽扰动改变了部分事件工期和命中率，但没有改变最终选中的切图、分核和核内顺序。失败候选不计入收益均值。
\paragraph{参数作用的直接判断。}容量减半使六图平均工期增加$0.2701\%$，容量加倍使其减少$0.0517\%$；变化集中在容量边界附近的图\Case{062}，其工期分别增加$1.6206\%$和减少$0.3103\%$。带宽减半和加倍的平均工期变化为$-0.0012\%$和$+0.0089\%$，方向未呈单调关系。四种扰动下重优化均回到同一份固定计划，说明在本组数据中，L2参数首先改变读请求的服务时序，尚未形成需要改变切图边界、核心归属或核内顺序的稳定收益。
\begin{table}[H]\centering\small\caption{六图四种Cache扰动的绝对工期均值与逐图命中率（固定计划与重优化）}\label{tab:cache-sensitivity}
\begin{tabular}{lrrrrr}\toprule
配置 & \shortstack{固定绝对工期均值\\/cycle} & \shortstack{重优化绝对工期均值\\/cycle} & \shortstack{逐图命中率\\均值} & \shortstack{逐图相对默认\\变化均值} & 失败候选\\\midrule
容量减半 & 4343744.7 & 4343744.7 & 10.71\% & +0.2701\% & 3\\
容量加倍 & 4338303.3 & 4338303.3 & 14.73\% & -0.0517\% & 3\\
带宽减半 & 4339193.8 & 4339193.8 & 14.63\% & -0.0012\% & 3\\
带宽加倍 & 4339315.0 & 4339315.0 & 14.63\% & +0.0089\% & 3\\\bottomrule
\end{tabular}\end{table}
表中两列工期均值分别对六张图的绝对周期取算术平均；“逐图相对默认变化均值”先按每张图的$100(F_{\rm variant}/F_{\rm default}-1)$计算，再在六图间等权平均。逐图比值的均值不能由两列绝对工期均值直接相除得到。
\samplefigure{sensitivity_makespan}{六图固定计划在四种Cache扰动下的平均工期变化}{fig:cache-sensitivity-time}
\samplefigure{sensitivity_hit}{六图固定计划在四种Cache扰动下的逐图命中率均值}{fig:cache-sensitivity-hit}

表中“命中率均值”先逐图计算再等权平均；若汇总六图的命中与访问字节，加权率依次为0.7876\%、1.6757\%、1.5597\%、1.5597\%。表\ref{tab:c-sensitivity-cases}显示工期响应集中于少数图：容量减半的平均工期增加0.2701\%，其中\Case{062}增加1.6206\%；容量加倍平均减少0.0517\%，仅该图出现0.3103\%的降幅。带宽扰动在多数图上不改变工期，少数图出现幅度很小的反向变化。六图固定计划的参数响应因而表现为稀疏、逐例不完全单调。
这组结果显示，容量扰动对命中和尾部工期的影响集中在驻留区间接近容量边界的图；带宽在本组固定计划上的平均变化仅为$-0.0012\%$和$+0.0089\%$，且方向不单调。四种扰动下的有限重优化均保留原计划，结合新增的领读核与错峰候选，本文把这一结果解释为“在当前六图和候选预算内未观察到计划改变”，而不外推为L2参数对切图边界没有作用。
表\ref{tab:c-sensitivity-cases}逐图给出$100(F_{\rm variant}/F_{\rm default}-1)$，正值表示变慢。每行四次比较使用同一默认五核计划。容量减半与加倍时，只有\Case{062}的工期变化；带宽减半影响两图，带宽加倍影响三图。逐图列出零值，使均值的来源与参数作用范围均可直接查验；带宽扰动出现的微小反向变化保留为事件级观察，不把它归因于某一条未展开的关键链。
\begin{table}[H]\centering\small\caption{六张代表图固定计划的Cache参数逐例响应（\%）}\label{tab:c-sensitivity-cases}
\begin{tabular}{lrrrrr}\toprule
图 & 默认工期/cycle & 容量减半 & 容量加倍 & 带宽减半 & 带宽加倍\\\midrule
\Case{012} & 13406 & 0.0000 & 0.0000 & 0.0000 & 0.0000\\
\Case{016} & 7740666 & 0.0000 & 0.0000 & 0.0000 & +0.0006\\
\Case{051} & 627502 & 0.0000 & 0.0000 & -0.0209 & +0.0105\\
\Case{062} & 1690862 & +1.6206 & -0.3103 & +0.0135 & +0.0420\\
\Case{067} & 15931687 & 0.0000 & 0.0000 & 0.0000 & 0.0000\\
\Case{081} & 30943 & 0.0000 & 0.0000 & 0.0000 & 0.0000\\\bottomrule
\end{tabular}\end{table}

命中率与工期并不同步。\Case{051}容量减半后字节命中率从42.39\%降至29.35\%，工期仍为627502 cycle；改变Cache读带宽时，命中率均为42.39\%，工期分别比默认短0.0209\%和长0.0105\%。\Case{062}容量减半使工期增加27402 cycle，容量加倍使工期减少5246 cycle，对应命中率为1.43\%和4.97\%。这些结果说明，命中增加能否缩短尾部完成时间，还取决于读事件所处执行路径及共享池的服务顺序。

\paragraph{灵敏度分析小结。}
问题三得到100张图逐核的C场景计划与工期，以及500组同计划B/C配对。五核固定计划中54图因Cache缩短工期，46图工期不变；进一步调整C计划仅在7图上取得改善。六图参数扰动显示容量和带宽的工期影响集中于少数图。模型将发射查询、完成插入、共享服务和最终完成时间串成同一条求解链；轻量层用收缩图回放查询相对时序，正式工期、命中和淘汰字段仍由有限候选的官方评价结果给出。
